% ============================================================
% 1.3 Fundamentación Teórica
% Archivo: tex/Theoretical-foundation/fundamento.tex
% Bibliografía: bib/fundamentacion_teorica.bib (+ references.bib compartido)
% Se conservan solo los conceptos empleados en el desarrollo del trabajo.
% ============================================================

\subsection{Fundamentación teórica}

Los conceptos relacionados con las variables del tema se organizan en cuatro
bloques: el flujo de información clínico-administrativo y su representación, los
sistemas de información en salud y su arquitectura de software, los agentes
conversacionales para la comunicación con el paciente y el marco normativo
ecuatoriano de la información clínica.

% -----------------------------------------------------------
\subsubsection{Gestión del flujo de información clínico-administrativa}
% -----------------------------------------------------------

El flujo de información clínico-administrativa comprende las actividades,
decisiones y transferencias de datos que ocurren desde el primer contacto del
paciente con el establecimiento hasta el cierre administrativo de su atención. Abarca
tanto los datos clínicos (diagnóstico, tratamiento, odontograma) como los de gestión
(agenda, cobro, seguimiento), y su representación formal hace visibles los pasos,
actores y dependencias que en la práctica permanecen implícitos.

\paragraph{Modelado de procesos de negocio (BPMN).}

El \textit{Business Process Model and Notation} (BPMN) es un estándar de notación
gráfica para representar procesos de negocio, mantenido por el \textit{Object
Management Group} (OMG) y publicado como ISO/IEC 19510. Su notación resulta
comprensible tanto para los actores del negocio como para los desarrolladores, de
modo que un mismo diagrama sirve de puente entre el proceso y su implementación en
software \cite{omg2024bpmn}. En el sector salud se ha empleado para documentar
procesos clínicos con roles simultáneos y alta variabilidad \cite{pufahl2022bpmn},
aunque el rigor con que se elaboran estos mapas de procesos varía considerablemente
entre estudios \cite{antonacci2021process}.

\paragraph{Estudios de tiempos y movimientos.}

Los estudios de tiempos y movimientos (\textit{Time and Motion Studies}, TMS)
provienen de la ingeniería industrial y se emplean en la investigación biomédica para
medir la duración de los eventos de un proceso asistencial \cite{cureus2022purpose}.
El término agrupa variantes que van desde la observación continua por un observador
externo hasta el análisis de registros del sistema; su único elemento común es la
captura de la duración de uno o más eventos \cite{Lopetegui2014TMS}. En el registro
dental electrónico, esta metodología ha permitido identificar interrupciones del flujo
de trabajo asociadas al registro de cargos, la revisión de radiografías y el uso de la
agenda \cite{schwei2016exploring}.

% -----------------------------------------------------------
\subsubsection{Sistemas de información en salud}
% -----------------------------------------------------------

Un sistema de información en salud es el conjunto organizado de datos, procesos y
tecnología que produce información para la toma de decisiones en una organización
sanitaria \cite{sciencedirect2024his}. Este marco orienta tanto la selección de sus
componentes funcionales como la evaluación de su éxito frente a los objetivos de la
organización y las exigencias regulatorias \cite{ovreiu2024empowering}.

\paragraph{Historia clínica electrónica (EHR).}

La historia clínica electrónica (\textit{Electronic Health Record}, EHR) reemplaza el
registro en papel por un repositorio digital estructurado de la información clínica
del paciente. Su adopción se asocia con mejoras en la eficiencia clínica y en la
gestión de datos, aunque persisten barreras de costo, seguridad y resistencia del
usuario \cite{sravya2024electronic}. En odontología, la falta de estándares comunes
para su implementación dificulta la comunicación entre plataformas y la continuidad
del cuidado \cite{gurupur2024standards}.

\paragraph{Metodologías de desarrollo de software.}

Las metodologías ágiles organizan el desarrollo en entregas incrementales con
retroalimentación frecuente \cite{sommerville2016software}. Scrum estructura el
trabajo en iteraciones de duración fija, con roles definidos para el equipo
\cite{schwaber2020scrum}. \textit{Extreme Programming} (XP) establece prácticas
técnicas, como la programación en pareja y el desarrollo guiado por pruebas
\cite{beck2004extreme}. Kanban visualiza el flujo de trabajo en un tablero y
limita el trabajo en curso \cite{anderson2010kanban}. El desarrollo rápido de
aplicaciones (RAD) organiza el proyecto en cuatro fases: planificación de
requerimientos, diseño con el usuario mediante prototipos, construcción y
transición \cite{martin1991rapid}.

\paragraph{Arquitectura de software: backend, frontend y API REST.}

Una interfaz de programación de aplicaciones (API) permite que una aplicación acceda
a los recursos de otra. El estilo REST (\textit{Representational State Transfer}), el
de mayor adopción en la web, se caracteriza por la independencia entre cliente y
servidor y por la identificación uniforme de los recursos \cite{ibm2024restapi}. El
\textit{backend} concentra la lógica de negocio y el acceso a datos, y el
\textit{frontend} presenta la información al usuario consumiendo esas interfaces.

\paragraph{Modelado de la arquitectura: modelo C4 y UML.}

El modelo C4 describe una arquitectura mediante cuatro niveles de detalle
progresivo: contexto, contenedores, componentes y código \cite{brown2018c4}. UML
complementa estas vistas con diagramas de secuencia, de estados y de despliegue,
que representan el comportamiento y la distribución física del sistema
\cite{omg2017uml}.

\paragraph{Microservicios y monolito modular.}

La arquitectura de microservicios estructura una aplicación como un conjunto de
servicios pequeños, desplegables de forma independiente y comunicados mediante API
REST, cada uno responsable de una capacidad de negocio \cite{ibm2026microservices}.
El monolito modular, en cambio, se despliega como una única unidad, pero divide su
código en módulos con interfaces públicas y sin dependencias internas cruzadas; así
conserva la separación de responsabilidades sin el costo operativo de distribuir y
coordinar servicios independientes \cite{richards2020fundamentals}.

\paragraph{Arquitectura de puertos y adaptadores.}

La arquitectura de puertos y adaptadores, o hexagonal, aísla la lógica de una
aplicación de sus dispositivos externos: la lógica define puertos, es decir,
interfaces, y cada tecnología externa se conecta mediante un adaptador que puede
sustituirse sin modificar la lógica \cite{cockburn2005hexagonal}.

\paragraph{Contenedorización y despliegue: Docker.}

Un contenedor es una unidad estándar de software que empaqueta el código de una
aplicación con todas sus dependencias, de modo que se ejecuta de forma consistente en
cualquier entorno. Docker es la plataforma de contenedores de mayor adopción en la
industria \cite{docker2024containers}.

\paragraph{Autenticación mediante JSON Web Token (JWT).}

Un \textit{JSON Web Token} (JWT) es un estándar abierto (RFC~7519) que transmite
información entre partes como un objeto JSON firmado digitalmente, lo que permite
verificar la identidad del emisor y la integridad del contenido sin consultar una
sesión en el servidor \cite{auth0jwt2024introduction}.

\paragraph{Diseño guiado por atributos (ADD).}

El método \textit{Attribute-Driven Design} (ADD), propuesto por el
\textit{Software Engineering Institute} (SEI), define la arquitectura a partir de
sus factores de diseño (\textit{architectural drivers}): los requerimientos
funcionales prioritarios, los atributos de calidad y las restricciones del sistema.
Estos factores guían las decisiones de diseño en cada iteración del método
\cite{cervantes2016designing}.

% -----------------------------------------------------------
\subsubsection{Agentes conversacionales para la gestión de citas}
% -----------------------------------------------------------

\paragraph{Agentes conversacionales y procesamiento de lenguaje natural.}

Un agente conversacional, o \textit{chatbot}, es un programa capaz de comunicarse con
personas en lenguaje natural, entender sus expresiones y generar respuestas
pertinentes; sus arquitecturas han evolucionado desde reglas escritas manualmente
hacia métodos estadísticos y de aprendizaje profundo \cite{csaky2019chatbot}. Para
interpretar el texto libre se apoyan en el procesamiento de lenguaje natural
(\textit{Natural Language Processing}, NLP), subcampo de la inteligencia artificial
que permite reconocer la intención de un mensaje y extraer sus datos
\cite{ibm2024nlp}.

\paragraph{Modelos de lenguaje grande (LLM).}

Un modelo de lenguaje grande (\textit{Large Language Model}, LLM) es un modelo de
aprendizaje profundo, con miles de millones de parámetros, entrenado sobre grandes
volúmenes de texto, capaz de comprender y generar lenguaje natural y de resolver
tareas nuevas a partir de instrucciones, sin un entrenamiento específico
\cite{zhao2023llmsurvey}. Puede emplearse para reconocer la intención de un mensaje
y extraer sus datos en un formato estructurado; sus respuestas, sin embargo, no son
deterministas y pueden contener información inexacta \cite{zhao2023llmsurvey}.

\paragraph{WhatsApp Business Platform.}

La \textit{WhatsApp Business Platform}, operada por Meta, permite a una organización
enviar y recibir mensajes de forma programática mediante su Cloud API. La integración
con sistemas externos se realiza mediante \textit{webhooks}: cuando ocurre un evento,
como la recepción de un mensaje, la plataforma envía los datos a una URL registrada,
que debe confirmar su recepción con una respuesta HTTP \cite{meta2026whatsapp}.

% -----------------------------------------------------------
\subsubsection{Marco normativo de la información clínica en Ecuador}
% -----------------------------------------------------------

\paragraph{Historia clínica y confidencialidad de la información en salud.}

El Acuerdo Ministerial 5216-A del Ministerio de Salud Pública del Ecuador establece
el marco obligatorio para el manejo de información confidencial en el Sistema
Nacional de Salud. Dispone que la gestión de las historias clínicas, en formato
físico o digital, garantice su archivo, custodia, confidencialidad y la trazabilidad
del uso de la información \cite{msp5216a}.

\paragraph{Registro clínico odontológico y codificación de diagnósticos.}

La historia clínica única odontológica se registra en Ecuador mediante un formulario
estandarizado del Sistema Nacional de Salud, el Formulario 033 \cite{msp5216a}. Su
completitud se evalúa con indicadores que combinan datos administrativos y clínicos,
entre ellos el índice CPO-D (dientes permanentes cariados, perdidos y obturados),
estandarizado por la Organización Mundial de la Salud para cuantificar la caries
dental \cite{who2013oralhealth}. Los diagnósticos se codifican con la Clasificación
Internacional de Enfermedades, décima revisión (CIE-10), también de la Organización
Mundial de la Salud, que asigna a cada enfermedad un código alfanumérico; las
afecciones bucales y dentales se agrupan en las categorías K00 a K14
\cite{who2019icd10}.

\paragraph{Protección de datos personales en el ámbito sanitario.}

La Ley Orgánica de Protección de Datos Personales del Ecuador regula el tratamiento
de los datos personales y considera los datos de salud como una categoría especial
\cite{lopdp2021}. Su cumplimiento en los sistemas de historia clínica electrónica del
país presenta brechas entre lo dispuesto por la norma y su implementación técnica
\cite{martinez2022proteccion}.
